\documentclass[11pt]{article}
\usepackage[utf8]{inputenc}
\usepackage[T1]{fontenc}
\usepackage[spanish]{babel}
\usepackage[margin=2.5cm]{geometry}
\usepackage{enumitem}
\usepackage{hyperref}

\begin{document}

\begin{center}
    {\large \textbf{Carrera de Ingeniería en Software}}\\[0.3cm]

    {\large \textbf{Universidad Técnica Estatal de Quevedo}}\\[0.2cm]
    {\large Facultad de Ciencias de la Computación y Diseño Digital}\\[0.3cm]

    {\large {Sistemas Distribuidos}}\\[0.2cm]
    {\large FORM-FB-01 - Fichas bibliográficas}\\[0.5cm]

    {\large {Giraldo Cagua Diego José}}
\end{center}

\vspace{1cm}


\section*{Ficha 1}

\subsection*{Título de la fuente}
\textit{A Survey on Microservices Architecture: Principles, Patterns and Migration Challenges}

\subsection*{Referencia IEEE}
V. Velepucha and P. Flores, ``A Survey on Microservices Architecture: Principles, Patterns and Migration Challenges,'' \textit{IEEE Access}, vol. 11, pp. 88339--88358, 2023, doi: \href{https://doi.org/10.1109/ACCESS.2023.3305687}{10.1109/ACCESS.2023.3305687}.

\subsection*{Resumen}
La fuente aborda la falta de una visión integral sobre los principios, patrones y retos asociados al escoger una arquitecturas de microservicios. Mediante una revisión, examina conceptos fundamentales, estrategias de descomposición de sistemas monolíticos, comunicación entre servicios, gestión de datos, despliegue y observabilidad. Entre sus resultados, identifica que la migración hacia microservicios enfrenta retos relacionados con la granularidad de los servicios, la consistencia de datos distribuidos, la seguridad y la complejidad operativa. Una idea clave es que la arquitectura de microservicios no es adecuada para todos los contextos. Su adopción debe justificarse según los requisitos de escalabilidad, autonomía de equipos y evolución del sistema. Esta conclusión aparece en la sección de discusión y conclusiones del artículo, particularmente en la síntesis de desafíos de migración.

\subsection*{Dato concreto}
El artículo señala que la identificación de límites adecuados entre microservicios es uno de los principales desafíos de la migración, ya que una granularidad incorrecta puede generar acoplamiento excesivo o servicios demasiado pequeños. Este planteamiento se desarrolla en la sección dedicada a los retos y principios de la arquitectura de microservicios, dentro del cuerpo del artículo.

\subsection*{Nota personal}
Este estudio ofrece un marco conceptual útil para comprender que la transición hacia una arquitectura de microservicios no debe considerarse simplemente como una decisión técnica; por el contrario, debe abordarse tanto desde una perspectiva arquitectónica como organizativa. Las ideas aquí expuestas subrayan que definir los límites adecuados de los servicios es fundamental para evitar problemas como el acoplamiento excesivo entre componentes, la complejidad operativa y las dificultades en el mantenimiento del sistema. Asimismo, el documento destaca que cualquier decisión de abandonar una arquitectura monolítica debe evaluarse en función de los requisitos específicos del negocio, las necesidades de escalabilidad previstas y la preparación del equipo para gestionar sistemas distribuidos.


\section*{Ficha 2}

\subsection*{Título de la fuente}
\textit{Balancing technology heterogeneity in microservice architectures}

\subsection*{Referencia IEEE}
G.-D. Schwarz, P. Heltweg, and D. Riehle, ``Balancing technology heterogeneity in microservice architectures,'' \textit{Empirical Software Engineering}, vol. 30, art. no. 127, 2025, doi: \href{https://doi.org/10.1007/s10664-025-10684-4}{10.1007/s10664-025-10684-4}.

\subsection*{Resumen}
Esta investigación analiza cómo la diversidad tecnológica influye en el diseño y desarrollo de arquitecturas de microservicios. Los autores destacan los desafíos que conlleva decidir cuándo resulta apropiado emplear distintas tecnologías en diversos servicios, subrayando la necesidad de encontrar un equilibrio entre los beneficios como la especialización técnica y la autonomía y los costos, entre los que se incluyen una mayor complejidad en la integración, el mantenimiento y la gestión. Con este fin, se llevó a cabo un estudio cualitativo que consistió en entrevistas a 15 expertos en desarrollo de software y arquitectura de sistemas. La investigación propone un marco para evaluar la diversidad tecnológica en los microservicios, clasificándola como relevante, aceptable o problemática. Un aspecto fundamental es que dicha diversidad debe estar justificada por requisitos funcionales u organizativos, o bien por especificaciones de rendimiento; por tanto, no debe adoptarse de manera indiscriminada. Esta conclusión se refleja tanto en las hipótesis como en el análisis de las implicaciones prácticas que se presentan en el trabajo.

\subsection*{Dato concreto}
El estudio se basó en entrevistas con 15 expertos, mediante una encuesta cualitativa orientada a construir una teoría sobre la gestión de la heterogeneidad tecnológica en arquitecturas de microservicios. Este dato aparece en la sección metodológica del artículo.

\subsection*{Nota personal}
Este recurso ofrece una perspectiva sobre el debate en torno a si la diversidad de tecnologías empleadas en los microservicios debe gestionarse y documentarse de manera transparente. El informe respalda la idea de que optar por una variedad de lenguajes de programación, marcos de trabajo o bases de datos puede resultar beneficioso para satisfacer requisitos específicos del mundo real, a pesar de que esto incremente la complejidad en términos de integración, despliegue y mantenimiento. En consecuencia, el recurso subraya la recomendación de definir tanto los requisitos técnicos como los organizativos antes de introducir nuevas tecnologías en una arquitectura distribuida.

\section*{Ficha 3}

\subsection*{Título de la fuente}
\textit{Full stack optimization of microservice architecture: systematic review and research opportunity}

\subsection*{Referencia IEEE}
M. Yang, X. Wang, B. Cai, et al., ``Full stack optimization of microservice architecture: systematic review and research opportunity,'' \textit{Cluster Computing}, vol. 28, art. no. 1005, 2025, doi: \href{https://doi.org/10.1007/s10586-025-05690-6}{10.1007/s10586-025-05690-6}.

\subsection*{Resumen}
Este recurso se centra en la optimización de la arquitectura de microservicios, destacando que el rendimiento y la eficacia generales dependen de las decisiones tomadas en diversos niveles del sistema. Un problema importante es la falta de cohesión en la investigación existente y gran parte del trabajo actual se concentra en optimizar componentes individuales como contenedores, comunicación, bases de datos o el escalado del sistema sin analizar las interdependencias entre ellos. Para abordar esta limitación, los autores presentan una revisión sistemática de la literatura sobre la optimización de microservicios a través de distintas capas tecnológicas. Los hallazgos permiten clasificar las estrategias de optimización e identificar nuevas líneas de investigación, especialmente en lo que respecta a la coordinación entre la infraestructura, la orquestación, la comunicación y las propias aplicaciones. Una conclusión clave es la necesidad de adoptar una perspectiva integral (de extremo a extremo), ya que la mejora de partes específicas del entorno no se traduce necesariamente en un mejor rendimiento global del sistema. Esta perspectiva surge de un análisis sobre futuras líneas de investigación y de una síntesis de los resultados obtenidos a partir de la revisión bibliográfica.

\subsection*{Dato concreto}
El artículo plantea que la optimización de microservicios debe abordarse desde una perspectiva de pila completa o \textit{full stack}, integrando decisiones de aplicación, comunicación, contenedores, orquestación e infraestructura. Esta definición y planteamiento aparecen en la introducción y en la sección de discusión del artículo.

\subsection*{Nota personal}
Este estudio ofrece un enfoque integral para analizar el rendimiento de sistemas basados en microservicios. El autor señala que optimizar un único componente como la contenedorización o la comunicación entre servicios resulta insuficiente si no se tienen en cuenta, además, las interacciones con la infraestructura, los sistemas de orquestación y las capas de aplicación. En consecuencia, el estudio subraya la importancia de considerar la escalabilidad, la tolerancia a fallos y los costes operativos como parte de una evaluación del rendimiento del sistema.


\section*{Ficha 4}

\subsection*{Título de la fuente}
\textit{Exploring microservice ownership and organizational coupling in open-source projects: an empirical study}

\subsection*{Referencia IEEE}
X. Li, D. A. d’Aragona, T. Cerny, et al., ``Exploring microservice ownership and organizational coupling in open-source projects: an empirical study,'' \textit{Computing}, vol. 107, art. no. 102, 2025, doi: \href{https://doi.org/10.1007/s00607-025-01454-7}{10.1007/s00607-025-01454-7}.

\subsection*{Resumen}
Esta investigación analiza la relación entre la asignación de responsabilidades en el contexto de los microservicios y el acoplamiento organizativo en proyectos de software de código abierto. Los investigadores señalan la falta de evidencia sobre cómo la asignación de equipos a servicios específicos influye en la coordinación de actividades, las interdependencias entre componentes y el proceso de desarrollo de software. Para abordar esta cuestión, el estudio lleva a cabo un análisis empírico de diversos proyectos de código abierto, examinando la asignación de responsabilidades en los microservicios, los cambios implementados por los distintos equipos y las relaciones o acoplamiento entre los componentes del sistema. Los resultados indican que compartir o distribuir responsabilidades entre los servicios puede derivar en un acoplamiento organizativo excesivo, lo cual dificulta la coordinación entre los equipos. Una conclusión fundamental es que unos límites de responsabilidad claramente definidos fomentan la autonomía de los equipos y reducen la necesidad de una comunicación innecesaria y esta observación se presenta como parte de la discusión del estudio y de sus correspondientes implicaciones prácticas.

\subsection*{Dato concreto}
El artículo concluye que una propiedad clara de los microservicios contribuye a reducir el acoplamiento organizacional, pues delimita las responsabilidades de los equipos y facilita la evolución independiente de los servicios. Esta conclusión aparece en la sección de discusión e implicaciones del estudio empírico.

\subsection*{Nota personal}
Este estudio ofrece perspectivas sobre la dimensión organizativa de una arquitectura de microservicios, subrayando que la partición de los servicios técnicos debe ir acompañada de la asignación de áreas de responsabilidad claras a grupos específicos. El establecimiento de áreas de responsabilidad bien definidas refuerza la autonomía operativa, reduce las dependencias entre componentes y permite el desarrollo independiente de cada servicio. Por consiguiente, el éxito de una arquitectura de microservicios no depende únicamente de su diseño técnico, sino también de su estructura organizativa y de sus principios de gobernanza.


\section*{Ficha 5}

\subsection*{Título de la fuente}
\textit{A Taxonomy of Integration-Relevant Faults for Microservice Testing}

\subsection*{Referencia IEEE}
L. Gregor, A. Hentschel, L. Kastner, and A. Pretschner, ``A Taxonomy of Integration-Relevant Faults for Microservice Testing,'' in \textit{2025 IEEE Conference on Software Testing, Verification and Validation (ICST)}, 2025, doi: \href{https://doi.org/10.1109/ICST62969.2025.10989000}{10.1109/ICST62969.2025.10989000}.

\subsection*{Resumen}
Este recurso aborda los desafíos que conlleva la identificación y clasificación de defectos que afectan a la interoperabilidad de los microservicios. Un aspecto fundamental es que las pruebas tradicionales centradas en componentes individuales suelen ser insuficientes para detectar errores de interoperabilidad entre servicios, tales como fallos en la comunicación, inconsistencias en los datos o problemas de dependencia en las interfaces. Para resolver esto, los autores han desarrollado un sistema de clasificación de defectos de interoperabilidad en microservicios, basado en el análisis de los tipos de defectos, los mecanismos de comunicación y los modos de fallo. El sistema de clasificación resultante pretende servir de guía para la selección y el diseño de pruebas de interoperabilidad.

\subsection*{Dato concreto}
El artículo propone una taxonomía específica de fallos relevantes para la integración de microservicios, orientada a mejorar la selección y el diseño de pruebas de integración. Esta contribución se presenta en la sección de propuesta de la taxonomía y en la discusión de resultados del artículo.

\subsection*{Nota personal}
Este estudio presenta un enfoque centrado en la calidad y la fiabilidad de los sistemas basados en una arquitectura de microservicios. El autor explica que las pruebas unitarias, por sí solas, resultan insuficientes para garantizar el correcto funcionamiento de una arquitectura distribuida, ya que muchos errores surgen de la interacción entre los servicios. La clasificación propuesta sirve como marco de referencia para identificar errores de comunicación, problemas de interfaz e inconsistencias en los datos. En consecuencia, el autor aboga por incorporar pruebas de integración, pruebas de consenso y pruebas de resiliencia en el proceso de desarrollo de microservicios.

\newpage

\section*{Declaración de uso de inteligencia artificial}

Para la elaboración de este trabajo se utilizó la herramienta de inteligencia artificial {Perplexity}.

La herramienta fue empleada para:

\begin{itemize}
    \item Localizar y proponer fuentes académicas recientes sobre sistemas distribuidos y arquitecturas de microservicios.
    \item Redactar las fichas bibliográficas en formato LaTeX.
\end{itemize}

El uso de la inteligencia artificial se realizó como apoyo en la búsqueda, organización y redacción preliminar del material. La información obtenida fue verificada mediante la revisión de los DOI, títulos, autores, años de publicación, editoriales y enlaces oficiales de las fuentes académicas. Asimismo, se comprobó que cada fuente correspondiera a una publicación indexada, un artículo de revista o un artículo de congreso, conforme a los requisitos establecidos para el trabajo.

\end{document}
